\chapter{Introduction}

\section{Background of the Study}
In modern capital markets, equity exchanges such as the National Stock Exchange of India (NSE) and Bombay Stock Exchange (BSE) serve as vital engines of economic wealth creation and corporate capital formation. In recent years, democratization of low-cost mobile brokerage applications, simplified KYC processes, and social media dissemination have brought tens of millions of first-time retail investors into active equity and derivatives trading.

However, rapid retail participation has concurrently heightened vulnerability to organized market manipulation. Foremost among these illicit activities is the classic \textbf{Pump-and-Dump (P\&D)} scheme. In a coordinated pump-and-dump operation, colluding market syndicates, unscrupulous operators, and unregulated social media influencers artificially inflate the trading volume and market price of small-cap, micro-cap, or illiquid penny stocks. This is executed through circular wash sales, aggressive order layering, and coordinated promotion across messaging platforms such as Telegram, WhatsApp, and YouTube channels. 

Once unsuspecting retail investors rush to purchase the asset under Fear-Of-Missing-Out (FOMO), the manipulators abruptly liquidate their positions at the top, leaving retail participants trapped in illiquid lower-circuit freezes with severe capital destruction.

Simultaneously, retail derivative participants trading benchmark index options (such as NIFTY 50 and BANK NIFTY) suffer from severe information asymmetry. Unlike institutional proprietary desks that continuously monitor real-time Open Interest (OI) shifts, trapped option writer unwinding, and Implied Volatility (IV) regimes, retail traders predominantly rely on lagging indicators or ungrounded speculative tips. 

\textbf{SentinelMarket} is an advanced, real-time market surveillance terminal and institutional derivatives playbook engineered to bridge this informational divide. It unifies mathematical anomaly detection for cash equities with algorithmic options risk management, delivering actionable intelligence directly to traders through an ultra-dense, Bloomberg-inspired dark terminal.

\section{Problem Statement}
The research and engineering in this project addresses four interconnected market failure domains:

\begin{enumerate}
    \item \textbf{Lagging Regulatory Intervention:} Regulatory authorities (such as SEBI) primarily identify and penalize pump-and-dump operators through post-mortem forensic audits months or years after the fraudulent event has concluded and retail capital has been wiped out. Retail investors require proactive, \textit{in-flight ($t=0$)} algorithmic alerts during the accumulation and pump phases.
    \item \textbf{High False-Positive Anomaly Classification:} Conventional anomaly detection models frequently misinterpret legitimate corporate breakouts, quarterly earnings surprises, or broad institutional accumulation in liquid securities as manipulative pumps. A rigorous surveillance system must statistically decouple pure volume anomalies from fundamental order-book liquidity (such as daily turnover in Crores).
    \item \textbf{Information Asymmetry in Options Derivatives:} Retail option buyers systematically lose capital to time-decay ($\theta$) and volatility crush ($\nu$) because they lack tools to identify whether institutional writers are actively defending strike walls or experiencing catastrophic gamma panic squeezes.
    \item \textbf{Cluttered and Distracting Retail Interfaces:} Typical retail trading software is bloated with delayed news feeds, visual noise, and confusing charting overlays. Fast trading execution requires an institutional-grade, keyboard-driven, high-density terminal interface.
\end{enumerate}

\noindent \textbf{Formal Problem Statement:} To design, formulate, implement, and evaluate an automated, real-time market surveillance platform that detects quantitative equity volume and price anomalies on Indian exchanges, evaluates institutional derivatives open interest dynamics, and renders deterministic risk ratings in an ultra-low latency Bloomberg-grade dark terminal.

\section{Objectives of the Study}
The research objectives are organized into primary quantitative modeling goals, secondary derivatives strategy formulation goals, and software engineering deliverables:

\subsection{Primary Quantitative Surveillance Objectives}
\begin{enumerate}
    \item Formulate a rolling 20-day standardized Volume Z-score ($Z_V$) and median volume ratio pipeline to isolate statistically anomalous trading activity.
    \item Implement Wilder's 14-period Relative Strength Index (RSI) momentum tracking and Bollinger Band upper barrier breach metrics.
    \item Design a composite 0--100 Manipulation Risk Scoring model that classifies securities into deterministic traffic-light risk tiers (Clean Flow, Momentum, Volume Surge, Critical Pump).
\end{enumerate}

\subsection{Secondary Derivatives Playbook Objectives}
\begin{enumerate}
    \item Parse official live exchange derivatives telemetry to extract complete 31-strike NIFTY 50 and BANK NIFTY option books.
    \item Formulate an Intraday vs. Positional Open Interest (OI) battleground model to pinpoint institutional support and resistance barriers.
    \item Develop an automated Trapped Writer Gamma Panic detection algorithm to identify short squeeze breakouts.
    \item Formulate a Strike Spread Width filter ($\approx 1.65\%$ spot boundary) to mathematically differentiate trending directional markets from range-bound sideways chop.
\end{enumerate}

\subsection{Software Engineering \& Deployment Objectives}
\begin{enumerate}
    \item Construct a high-concurrency, asynchronous Python FastAPI backend service with sub-second response times and synthetic error handling.
    \item Develop a professional, high-density Bloomberg Dark Terminal user interface using Next.js 15, React 19, and Tailwind CSS v4 featuring global keyboard navigation hotkeys.
    \item Build automated test suites verifying mathematical scoring integrity, synthetic offline mock data, and API contract correctness.
\end{enumerate}

\section{Scope of the Study}
The functional scope of SentinelMarket encompasses:
\begin{itemize}
    \item Ingestion of live and historical tick-level equity data across the NIFTY 50 index universe and customizable user watchlists.
    \item Mathematical computation of volume Z-scores, price-change velocities, RSI momentum, and volatility envelopes without heavy offline machine learning black boxes.
    \item Real-time options chain parsing for active weekly and monthly index expiry contracts.
    \item Interactive, client-side filtering, live search with autocomplete, and keyboard-driven trading workflows.
\end{itemize}

\noindent \textbf{Delimitations:} The platform is designed as an algorithmic decision-support tool. It does not execute live monetary broker orders, manage trading capital, or provide personalized investment advisory services.

\section{Significance of the Study}
This research contributes directly to retail investor protection and financial market transparency. By open-sourcing an end-to-end, reproducible quantitative surveillance system, retail market participants gain access to algorithmic tripwires historically reserved for institutional hedge funds and exchange compliance departments. 

Furthermore, the project bridges the gap between academic anomaly detection models and practical, real-time user experience design, demonstrating that mathematically robust market monitoring can be delivered with minimal latency on commodity computing hardware.

\section{Research Questions}
The research questions mapped in Table \ref{tab:rq_mapping} guide the design, execution, and evaluation of SentinelMarket.

\begin{table}[H]
\centering
\caption{Research Questions and Evaluation Mapping}
\label{tab:rq_mapping}
\begin{tabularx}{\textwidth}{@{}lXl@{}}
\toprule
\textbf{ID} & \textbf{Research Question} & \textbf{Evaluation Status} \\
\midrule
RQ1 & Can a standardized Volume Z-Score ($Z_V$) reliably isolate suspicious liquidity surges prior to price collapse? & Validated via Historical Case Studies \\
RQ2 & Does combining RSI momentum divergence with volume ratios reduce false-positive manipulation flags in liquid blue-chip equities? & Empirically Verified in Screener Engine \\
RQ3 & Can intraday changes in Open Interest ($\Delta\text{OI}$) effectively detect trapped option writers before directional gamma breakouts? & Implemented in Live Options Playbook \\
RQ4 & What performance advantages does a zero-database, memory-cached architecture provide for sub-second real-time market surveillance? & Verified in Automated Testbed \\
RQ5 & Does an ultra-dense, monospace Bloomberg terminal design enhance cognitive clarity and execution speed for retail traders? & Implemented and UI Verified \\
\bottomrule
\end{tabularx}
\end{table}

\section{Organization of the Report}
The remainder of this report is structured as follows:
\begin{itemize}
    \item \textbf{Chapter 2 (Literature Review):} Synthesizes existing academic research on market manipulation, examines related machine learning and statistical methodologies, and establishes the formal research gap.
    \item \textbf{Chapter 3 (Methodology):} Details the mathematical formulations, quantitative risk metrics, derivatives strategy logic, architectural diagrams, and evaluation criteria.
    \item \textbf{Chapter 4 (Implementation):} Explains the engineering implementation across the FastAPI backend, PumpWatch mathematical engine, Next.js terminal frontend, and local execution pipelines.
    \item \textbf{Chapter 5 (Results and Discussions):} Presents experimental screening results, automated test suite verification, and comparative performance analyses.
    \item \textbf{Chapter 6 (Conclusion and Future Scope):} Summarizes project contributions, discusses operational limitations, and outlines future avenues for research.
    \item \textbf{Annexure \& References:} Contains code artifacts, setup procedures, and academic citations.
\end{itemize}
